\subsection{THETA3-A14. A prescribed tame primary family in one actual \texttt{TwDat}}
\label{A14}

This is a construction inside the fixed full host, rather than a change
of its CM field, ranks, or signatures.  Write $E/F$ for the host CM
extension, and $c$ for complex conjugation.  We retain the archive
hypotheses used in A13: $E/\Q$ is Galois and $E/F$ is unramified at every
finite place.  We also retain the entire \texttt{Pre}, the reference
character frame $f$, and the complex images of the two distinguished
coefficient subfields $K_s$ and $L_0$.  The extension of the larger
coefficient field $L$ below is exactly an \texttt{Enl} in the sense of
\texttt{CmQuadraticFormEqAlternatingForm.lean}, lines 71--74.

Here and below a primary character means a finite-order character of
prime-power order; it does not mean a character of residue-characteristic
power conductor.  In this construction the conductor is tame and is
prime to the primary prime.

\paragraph{Statement and the literal extra local condition.}
There are infinitely many odd rational primes $p$, avoiding any
prescribed finite set and $|\mu(E)|$, with the following property.
Let $\mathcal S$ be a finite collection of split places of $F$, lying
over rational primes that split completely in $E$ and avoid
\texttt{offSet} and $p$.  At each $v\in\mathcal S$ choose an ordering of
its two places in $E$.  Given finitely many systems
\[
 \lambda_j=\bigotimes_{v\in\mathcal S}\lambda_{j,v},
 \qquad 1\leq j\leq t,
 \tag{A14a}
\]
where $\lambda_{j,v}$ is any $p$-primary character of
$\kappa_v^\times$, there exist finite-order anticyclotomic Hecke
characters $\psi_j$ whose restriction at this split pair is
$\lambda_{j,v}(a/b)$.  They are unramified outside the indicated split
pairs, are tame everywhere, have $p$-power order, and are
\emph{trivial on all of $E_w^\times$} at every $w\mid p$.

After a finite normal coefficient-field enlargement, these characters
occur together in a single literal \texttt{TwDat}, with
\[
 \texttt{pl}=\{p\},\qquad
 \texttt{TwDat.set}(f)=
 \{1,\delta,\delta'\}\ \cup\ \{\delta\psi_j:1\leq j\leq t\}.
 \tag{A14b}
\]
Repeated characters are of course discarded in this set.  We choose
$p\ne p_{\delta},p_{\delta'}$.  If there are no prescribed characters,
the untwisted datum suffices.

The distinction between unramified and locally trivial is essential.
The actual definition is
\texttt{Formalism/LocalToGlobal/UnitaryNorms/CharHost.lean},
\texttt{IsPsi}, lines 77--84.  Besides order, anticyclotomicity,
tameness, exact ramification set, and split ramification, its
\texttt{atl} clause requires $\psi_w(x)=1$ for \emph{every}
$x\in E_w^\times$ when $w\mid p$.  The construction proves this clause,
not merely triviality on local units.

\paragraph{Proof of the character construction.}
Complex conjugation is central in $\operatorname{Gal}(E/\Q)$.
Choose $p$ unramified in $E$ with Frobenius $c$ and outside the indicated
finite set.  Chebotarev supplies infinitely many such primes; the precise
primary reference used here is Milne, \emph{Algebraic Number Theory},
version 3.08 (2020), Theorem 8.31 and Corollary 8.32, printed page 147.
Its restriction to $F$ is trivial, so $p$ splits completely in $F$;
at each resulting place, $E_w/F_v$ is the unramified quadratic
extension.  The finite exclusions include $2$, $|\mu(E)|$, the two
reference primes, and any already fixed bad primes.

Use the ratio convention at a split pair exactly as in A13.  Put
\[
 R_{\mathcal S}=\prod_{v\in\mathcal S}\kappa_v^\times,
 \qquad
 C_{\mathcal S}=
 \A_{E,f}^\times/
       \bigl(E^\times\A_{F,f}^\times U(\mathcal S)\bigr).
\]
Here $U(\mathcal S)$ consists of the principal local units at the two
places above each $v\in\mathcal S$ and all units elsewhere.  This is a
finite quotient of a ray class group.  A13 proves, including the
unramified valuation cancellation, that the kernel of
$R_{\mathcal S}\longrightarrow C_{\mathcal S}$ is the image of
$\mu(E)$ under the split ratios.  Each $\lambda_j$ has $p$-power
order.  Because $p\nmid |\mu(E)|$, its restriction to this kernel is
trivial.  Thus it defines a character on the image of
$R_{\mathcal S}$, and the elementary finite-abelian character extension
argument in A13 extends it to $C_{\mathcal S}$.

One can make this extension $p$-primary without changing the prescribed
local characters.  Write
$|C_{\mathcal S}|=p^a M$ with $(p,M)=1$, and choose an integer $e$ with
\[
 e\equiv1\pmod{p^a},\qquad e\equiv0\pmod M.
\]
Raising any extension to the $e$th power kills its prime-to-$p$ part
and acts identically on its prescribed $p$-primary restriction.  The
result is $\psi_j$.  Extend it trivially at infinity.  The defining
quotient gives triviality on $E^\times$ and on the finite ideles of
$F$; with trivial infinity type it is anticyclotomic on all ideles.
The subgroup $U(\mathcal S)$ gives tameness and the asserted conductor
support.  Let $Q_j$ be its exact finite ramification set, possibly
smaller than the prescribed collection of pairs.

At a place $w\mid p$, the character is unramified.  Since
$E_w/F_v$ is unramified, a uniformizer of $F_v$ is also one of $E_w$,
and therefore
\[
 E_w^\times=F_v^\times\,\OO_{E_w}^\times.
 \tag{A14c}
\]
Anticyclotomicity makes the local character trivial on the first
factor; unramifiedness makes it trivial on the second.  This proves
\texttt{atl}.  Every place in $Q_j$ is split, has prime absolute norm,
and lies above a rational prime outside \texttt{offSet}.  Such a prime
splits completely in $E$, so its ramification index is one.  These are
exactly the further clauses of \texttt{Chr.IsTwist}.  The defining
source is
\texttt{Development/AbelianVarietyTheory/Definitions/AutomorphicMembership/}
\texttt{IsRationalUnitaryDeterminantMapRationalUnitaryTo.lean},
lines 111--114.  We have proved all clauses of
\texttt{IsPsi} and \texttt{IsTwist}, with $Q_j$ the exact ramification
set.  In particular $Q_j$ contains no place above $p$.

There is no shortage of possible conductor primes: those splitting
completely in $E(\zeta_p)$ avoid any fixed finite set and satisfy
$q\equiv1\pmod p$.  Hence nontrivial tame $p$-primary characters occur
at infinitely many allowed primes.  This assertion does not prescribe
the power-residue values of a previously fixed finite list of matrix
entries.

\paragraph{Enough coefficient-place orbits, with a degree bound.}
The index $w$ of \texttt{TwDat} is a finite place of the coefficient
field, not a place of $E$.  First enlarge the normal coefficient field
$L$ to contain all finitely many character values and any already
specified algebraic constants.  For $t\geq1$ choose $r\geq0$ with
$2^r\geq t$.  There is a totally real multiquadratic extension
\[
 M=\Q(\sqrt{q_1},\ldots,\sqrt{q_r}),
 \qquad [M:\Q]=2^r,
 \tag{A14d}
\]
which is linearly disjoint from $L$, and in which $p$ splits
completely.  To construct it, choose distinct rational primes
$q_i\equiv1\pmod{4p}$, none dividing the discriminant of $L$.
The split-prime case of the same Chebotarev theorem applied to
$\Q(\zeta_{4p})$ supplies them.  Independence of their rational
squareclasses gives the stated degree.  Every nontrivial quadratic
subfield of $M$ is generated by the square root of a nonempty product
of the $q_i$ and is ramified at at least one of these primes.
It cannot be contained in $L$.  A nontrivial subfield of a
multiquadratic extension contains a quadratic subfield, so
$M\cap L=\Q$.  Also $q_i\equiv1\pmod p$, and Hensel's lemma at the
odd prime $p$ makes $q_i$ a square in $\Q_p$.  Thus $p$ splits
completely in $M$.

Replace $L$ by $LM$.  Each prime of $L$ above $p$ has $2^r$ primes
above it in $LM$: locally this is
$L_w\otimes_{\Q_p}(M\otimes_\Q\Q_p)
 \simeq L_w^{\,2^r}$.
Complex conjugation fixes $M$ pointwise, so it fixes each prime of
$M$ above $p$ and acts only on the $L$ component of these labels.
Consequently every conjugation-orbit of primes of $L$ above $p$
produces $2^r$ distinct conjugation-orbits in $LM$.  There are at
least $t$ such orbits.  The extra degree for this orbit-separation
step is
\[
 [LM:L]=2^{\lceil\log_2 t\rceil}<2t
 \quad(t>1),
 \qquad [LM:L]=1\quad(t=1).
 \tag{A14e}
\]
This bound concerns only the orbit-separation step.  It excludes
the earlier adjoining of character values and the later normalized
jet constants.

Choose $t$ distinct conjugation-orbits of coefficient primes above
$p$.  On the $j$th orbit set
$\psi(w)=\psi_j$, $k(w)=k_j$, and $Q(w)=Q_j$, where
$\psi_j^{p^{k_j}}=1$.  On every remaining coefficient prime put
$\psi(w)=1$, $k(w)=0$, and $Q(w)=\varnothing$.
Set \texttt{pl}$=\{p\}$.  The full definition
\texttt{CmQuadraticFormEqAlternatingForm.lean}, lines 79--89,
now checks clause by clause.  The three character conditions and
\texttt{Q\_l} were just proved; \texttt{off} holds by definition;
\texttt{conj} holds because the assignment is constant on each orbit.
At non-$p$ places the trivial character satisfies \texttt{IsPsi} for
their own residue characteristic.  At $p$ the prescribed characters
satisfy \texttt{IsPsi} by (A14c).  The formula for
\texttt{xiTw}, lines 91--95 of that source, gives (A14b), because
$p\ne p_\delta$ and all other assigned characters are trivial.

This datum survives every subsequent finite normal coefficient-field
enlargement.  Assign to a new coefficient prime the data of its
contraction to the old field.  Contraction preserves residue
characteristic and commutes with complex conjugation; lying over
ensures that every old character still occurs.  Thus all
\texttt{TwDat} clauses and its exact character range persist.
The A13 construction of \texttt{Base.Enl} keeps the two distinguished
subfields as their old images, so it also preserves the literal
\texttt{Base}, not just the list of fields.

\paragraph{An actual correct-range TC and, from an existing witness, ColW.}
For the TC assertion retain A13's diagonal-space and
\texttt{duStar} identities, the required reference-character values,
and its oriented \texttt{PreX}, valid selection, and valid
\texttt{AmbIn}.  A retained $K.\texttt{Over}\,\texttt{Pre}$
supplies the first two geometric identities and
\texttt{HasLetters}; otherwise state those conditions separately.
Apply A13 to the finite set (A14b), adjoining its finite
\texttt{Holds} constants at depth $2m$ and, if requested, the finite
normalized-jet generators for a fixed valid \texttt{AmbIn}.
The preceding pullback of \texttt{TwDat} preserves its range through
these extensions.  The exact source
\texttt{Development/AbelianVarietyTheory/Proofs/IwasawaDecomposition/}
\texttt{ExistsColumnInputAssembled.lean},
\texttt{isTwist\_of\_mem\_set}, lines 301--315, proves that every
member of this range is an allowed twist for
\[
 S=\{q:q\in\texttt{offSet},\ q\ne p_\delta,p_{\delta'}\}.
\]
Its proof includes both reference characters and their products with
the new characters.  The same file's
\texttt{exists\_colIn}, lines 195--289, then gives the correct-range
TC with the retained orientation, selection and valid ambient datum;
A13 supplies its canonical constants, \texttt{Laws} and
\texttt{Law0}; it supplies \texttt{Later} when
\texttt{HasLetters} holds.  Thus the new family is not merely
a collection of abstract Hecke characters.

There is also a full \texttt{ColW} statement with a concrete starting
hypothesis.  Suppose a witness \texttt{ColW} for this same
\texttt{Pre} and frame is already given, and the retained context
satisfies $K.\texttt{Over}\,\texttt{Pre}$.  Then the preceding construction
produces a new \texttt{ColW} with character range (A14b).
Here is the verification of the additional fields, which is needed
beyond the TC construction.  Orient its \texttt{PreX} using
\texttt{exists\_oriented}, lines 91--103 of
\texttt{ExistsColumnInputAssembled.lean}.  That explicit constructor
changes only the two \texttt{piG} choices and leaves
\texttt{piT} unchanged.  The two clauses of \texttt{TC.Pins}
in \texttt{Definitions/AutomorphicMembership/LatAtLattice.lean},
lines 125--130, involve only \texttt{Pre} and \texttt{piT}; hence they
are retained.  The equality \texttt{KsExact} is an equality of the
fixed complex field $K_s$ with the field generated by the unchanged
reference data, so it is retained as well.  The original
\texttt{ValLL0} asserts that those reference values lie in the fixed
complex image of $L_0$; it is likewise unchanged.

Finally \texttt{FieldOK} consists of the presence of the fixed
elliptic $j$-invariant and the principality of specified extended
fractional ideals.  This is its literal definition in
\texttt{Definitions/AutomorphicMembership/IsHalfPtIff.lean},
lines 31--41.  A generator of a principal ideal remains a generator
after coefficient extension: extending the equality
$I=x\OO_L$ gives $I\OO_{L'}=x\OO_{L'}$.  The defining spanning
generators of \texttt{ellLat} are the same elements of $E$, so its
new ideal is precisely that extension.  The $j$-invariant remains
in the larger field.  Thus \texttt{FieldOK} is preserved.
Combine these data with the new TC, its proven \texttt{Laws} and
\texttt{Law0}.  They are exactly all fields of \texttt{ColW},
\texttt{Definitions/AutomorphicMembership/PositionOfIntro.lean},
lines 23--36.  The retained $K.\texttt{Over}$ gives
\texttt{HasLetters}, and hence \texttt{Later}, as in A13.
This argument does not manufacture the original full host or
the starting \texttt{ColW} from a numerically matching generic CM torus.

\paragraph{Exactly which local character span has been obtained.}
At a chosen conductor prime write
$|\kappa_v^\times|=p^a h$, with $(p,h)=1$, and let
$H_v\subset\kappa_v^\times$ be its unique subgroup of order $h$.
The $p$-primary characters are precisely the character group of
$G_v=\kappa_v^\times/H_v$.  By finite Fourier inversion they span
all functions on $G_v$; the nontrivial ones span its mean-zero
subspace.  Explicitly, for distinct $z_0,z_1\in G_v$,
\[
 \mathbf1_{z=z_0}-\mathbf1_{z=z_1}
  =\frac1{|G_v|}\sum_{\chi\ne1}
       \bigl(\overline{\chi(z_0)}-\overline{\chi(z_1)}\bigr)\chi(z),
 \qquad
 \sum_{\chi\ne1}
   \left|\frac{\overline{\chi(z_0)}-\overline{\chi(z_1)}}{|G_v|}\right|
   \leq2.
 \tag{A14f}
\]
Every finite list of characters used in this identity can occur
together in the single datum just constructed.  On the literal
first-column \texttt{PhiQ} seed, this is a function of the
nonzero oriented wedge modulo $H_v$, retaining the full-rank support
and the right-$\mathrm{SL}$ covariance proved in A13.  It is not a
delta function on all matrices or even on all nonzero determinants.
At different primes the prescribed systems can be chosen independently,
so the corresponding tensor spans also occur in a common datum.
Actual global arrays still include their other local factors and
their class equivariance.  Their theta sections are obtained through
the literal finite Heisenberg transform, not by replacing those
factors with arbitrary Schwartz functions.

The new conclusion is the exact specialization of a prescribed finite
primary local family to one \texttt{TwDat}, and, with the stated
starting witness, to one \texttt{ColW}.  The remaining derivative
question is now a concrete finite matrix on this specified quotient
character space together with the incidence space of A13.  Arithmetic
inversion of such a matrix still requires a nonzero minor and the
section and tangent normalizations of A8 and A12.  In the requested
Grothendieck framework, this is a finite period comparison of the
kind isolated in \emph{cote44}, scans 7--8; it does not by itself
supply the cohomological correspondence or algebraic Lefschetz
inverse discussed in \emph{cote16}, scans 43--46.
